\chapter{岁计与人口：一座从来没有装满的城}

简十四是全书的枢纽：计口三十一万二千，实居十一万九千。

这一片之所以重要，不是因为它给出了一个数，而是因为它把两个本不该并列的数并列了。额定与实居之间差十二万一千，而写这片的人显然认为这个差额不需要解释。

\kao{一\quad 三十一万二千是什么数}

这是城的额定容量。简十说是「岁计口」的总数，简十二说是按此数配给的九年粮食。

但这个数在城建好之前就已经过时了。按《战争的雷霆》所复原的时序，城是为三十一万二千人设计并开工的，而完工的时候，这个文明只剩下十五万。

也就是说，这座城从落成那天起就没有可能装满。后来降到十二万，再降到十一万九千，只是把这个空置比例又往前推了一点。

\kao{二\quad 十一万九千是清点出来的}

简二十说「元年，会计。在籍十一万九千。可去者十一万九千」。

这两句其实是三句话：「在籍」是一个数，「可去者」是一个数，而这两个数一样大。第三句「留者无」才是结论。

在籍而不在城的人有没有，本卷没有说。简二十八说「外工三百二十人，不入装载之数。留」。若此片可信，则至少还有三百二十人不在这个数里。

\kao{三\quad 三万的差额}

从三十一万二千到十一万九千，中间掉了十九万三千。这个数太大，不能归给单一原因，故简十九把它拆成三项：兔乱、他故、逃亡。

三项合计二万二千。

还差十七万一千。

这十七万一千不是事故，是**两万年里自然发生的事**。一座两万年不生育的城，人口只减不增；每死一个人，在册人数就少一个，而新入者不足以补上——简十三说「新入者，先居外舍三月」，可见入城是特例，不是常态。

所以真正的账是这样：从三十一万二千降到十五万，是城的生命周期里自然发生的；从十五万降到十二万，是六批迁移的结果；从十二万降到十一万九千，是最后几年的零头。

简十九的三项只解释了其中的一小段。

\kao{四\quad 那八千}

上一章末尾提到的八千，是我在这一卷里唯一一处无法归位的数据。

简十九记事故二万二千，简二十记实存十一万九千，两者相加是十四万一千。而完工时的人数按《战争的雷霆》是十五万。

差八千。

这八千可能有三种去处：一是在简十九所记的三项之外，另有未记之事；二是简十四的「十一万九千」为取整，实际未必是这个数；三是十五万这个数本身的来源有问题，它出自《战争的雷霆》据另一批材料所作的复原，不是本卷的材料。

我倾向前者，但没有证据。这一条列在此处，不列入上编存疑，因为它的材料是清楚的，不清楚的是推论。

\kao{五}

最后说一句关于数字方法的话。

本卷的数字，凡能互相校验的，我一律校过。校不上的，写在这里，不写进正文。读者若认为某一处该有别的算法，尽管另算；原件都在附录一里列着，谁都可以重来。

这不是谦辞。一个只有三十七片残简的账，本来就凑不平。

\chapter{撤离：三因与六批}

简二十九、简三十两片，是全书信息密度最高的地方。前者给出六批的人数和总数，后者给出这六批之间的性质差别。

\kao{一\quad 六批}

\begin{center}
\small
\begin{tabular}{@{}llr@{}}
\toprule
批次 & 时间（据《战争的雷霆》） & 人数\\
\midrule
第一批 & 约前 421 年 & 6\,000\\
第二批 & 约前 380 年 & 8\,000\\
第三批 & 约前 340 年 & 10\,000\\
第四批 & 约前 300 年 & 13\,000\\
第五批 & 约前 260 年 & 18\,000\\
第六批 & 前 221 年 & 65\,000\\
\midrule
合计 & & 120\,000\\
\bottomrule
\end{tabular}
\end{center}

本卷只给出数字，不给年份。年份来自《战争的雷霆》所复原的编年。两书的数字在此处完全吻合，这是我把它们放在同一张表里的唯一理由。

\kao{二\quad 「岁计之内」与「岁计之外」}

简三十是全卷最要紧的一句话：「前五批，岁计之内。第六批，岁计之外。」

什么叫岁计之内，简十七给了例子：有定额，有赏罚，有退出机制。前五批是照这个办法走的：定人数，定次序，定配额，一批一批下去。

第六批不是。第六批六万五千人，在一年之内走完，占全部撤离人数的五成四。这一批没有配额，没有次序，只有「可去者十一万九千」。

\kao{三\quad 为什么走得这么急}

简二十七：「地热之衰，二百年间，百三十一分。始缓，终急。」

「始缓，终急」四字，是撤离近因的原件依据。热流下降不是匀速的，前期慢，后期快。到了某个点，维持气路与水路所需的能量超过可获得的总量，城就无法再维持。

《战争的雷霆》把这件事写成「纯粹工程问题，与道德判断无关」。本卷的材料支持这个说法：简二十一列了十九项退役措施，没有一项是出于安全考虑，全部是为了封存一座不再运转的设施。

\kao{四\quad 那三百二十人}

简二十八：「外工三百二十人，不入装载之数。留。」

「留」字之后有缺文。若此字属实，则撤离时有三百二十人被留下。

这与简二十「留者无」直接冲突。两片相隔不过八号，是同一批编目的材料。

我没有办法判定哪一片可信。若按简二十一「第九项：外舍，平」，垣外小舍在退役时被推平；若那三百二十人还在舍中，则他们被埋在那里。

这一条我不下结论。

\kao{五}

撤离之后，城的记录就断了。

简七十二片之后，全部是后世题记（卷四）。它们之中最早的一片说「后三百七十二年，有客至」。若自前 221 年计，那是公元 151 年。

三百七十二年里，这座城没有留下一个字。

\chapter{三事合考：射日、密葬、奔月}

这一章处理三件传世文献里的事。三件事表面上是三件，实际上指的是同一段历史的三个片段。

本卷没有直接材料。这一章的全部依据，是《战争的雷霆》第一章第四节、第三章第一节所引的记载，以及残简中若干间接痕迹。

\kao{一\quad 射日}

「十日并出，羿射九日」，历来有四说：天象、旱灾、历法、政治寓言。

本书取历法说，理由是简二十六的存活率表。

四代人，存活率从六一四降到四二，中间没有一次骤降。若是天象或旱灾，应当是突发的、成批的死亡；而这条曲线是连续的、逐代下滑的。它更像一种慢性病，或者一种长期不适应的环境。

若「十日」指十部历法，则「射九日」指废掉其中九部，只留一部。这与简十三「新入者，先居外舍三月」的制度精神一致：一切以统一为准。

要补充的是：后羿至此为止。他改了历，但据《战争的雷霆》所引，他一直没有发现那部改过的历还有别的用途。这件事在他死时仍然没有被发现。

\kao{二\quad 天文密葬法}

这个名字是月之民的，不是后羿的。

「密葬」二字的意思不是秘密埋葬，是「不留下可读的形」。一部历法写成这样，读它的人只能算出结果，不能反推出它算的是什么。

简三十五「□有十二，不可尽。尽者三，余九亡」，可能与这件事有关，但文中无一处可对上。我不据此立论。

\kao{三\quad 奔月}

「嫦娥奔月」在传世文献里有四个母题：服药、飞升、化为蟾蜍、独居月宫。

按《战争的雷霆》所复原，这四个母题对应一次事件：服药，被派下地，被某物攻击，然后被接回。

本卷对此只有一条间接材料：简三十一「门闭之日，无送者，无哭声。闭者七人，回首一次」。

「无送者」三字，若指的是撤离当日无人相送，则可与《战争的雷霆》所说「嫦娥没有随行」相印证。但本片无一处提到她。

我不把这两条对上。它们之间的关系，本卷的材料不足以判定。

\chapter{名号的层累}

这一章只做指认，不做发挥。

传世文献里的西王母、嫦娥、八意永琳、蓬莱山辉夜，在档案系统里对应的是同一批人。这一点必须说清楚，否则读者读别的卷时会对不上号。

\begin{center}
\small
\begin{tabular}{@{}lll@{}}
\toprule
文献所见 & 实际所指 & 出处\\
\midrule
嫦娥 & 驻地球主任 & 中原人所叫\\
西王母 & 同上一位的上级 & 中原人所叫\\
八意永琳 & 即西王母，后东渡日本 & 日本所叫\\
蓬莱山辉夜 & 东渡者之一 & 日本所叫\\
\bottomrule
\end{tabular}
\end{center}

这张表只解决「谁是谁」，不解决「名字怎么变的」。

「嫦娥」原作「姮娥」，避汉文帝刘恒讳而改；「天文密葬法」是技法之名，被后世误当作人名；「蓬莱」本为地名，后成为一部书的书名。这些变化各有各的路径，彼此并不相干，把它们统起来讲，会讲成一个并不存在的故事。

有一个尴尬必须处理。

这些名字在通俗读物里另有来历：一部作品，一群爱好者，一整套设定。读者若去查，查到的大概是那一片热闹，不是这一张表。

本书与那是两回事，不必互相印证。

\begin{flushright}
{\small 巴兰尼科夫}\quad{\footnotesize\color{pencil} 二〇二四年十一月，于莫斯科}
\end{flushright}
